
\section{Experiments}
\label{sec:experiments}

\subsection{Experimental Setup}
\label{sec:exp_setup}

\paragraph{Base model.}
We evaluate \rasp{} on DeepSeek-R1-Distill-Qwen-7B~\citep{deepseek2025r1}, a 7B-parameter reasoning model built on the Qwen2.5 architecture~\citep{qwen2.5} with 28 transformer layers.
The model uses Grouped Query Attention (GQA;~\citealp{ainslie2023gqa}) with 28 query heads and 4 key-value heads per layer (group size $G{=}7$), and an MLP intermediate dimension of 18{,}944.
As a reasoning-distilled model whose chain-of-thought capabilities originate from reinforcement learning, it provides a direct testbed for studying how structural pruning interacts with RLVR-based recovery.

\paragraph{Evaluation.}
We evaluate on GSM8K~\citep{cobbe2021gsm8k}, a benchmark of 1{,}319 grade-school math problems requiring multi-step arithmetic reasoning.
We report exact-match accuracy: the final numerical answer is extracted from the model's generation and compared against the ground truth.
For rapid iteration during development, we use a fixed 100-problem subset ($N{=}100$); final results use the full test set ($N{=}1{,}319$).
We define recovery $\Delta$ as the accuracy gain from post-prune to post-recovery, capturing how much reasoning ability the recovery pipeline restores independent of initial pruning damage.

\paragraph{Hardware.}
All experiments run on a single node with 2$\times$ NVIDIA RTX PRO 6000 GPUs (94\,GB VRAM each).
The two pruning variants (Standard and \rasp{}) are trained in parallel on separate GPUs to ensure identical wall-clock conditions.

\subsection{Pruning Configuration}
\label{sec:pruning_config}

\paragraph{Sparsity and granularity.}
We apply 20\% structured sparsity, removing entire attention GQA groups (each comprising 1 KV head and 7 query heads) and MLP channel groups jointly.
The MLP intermediate dimension (18{,}944) is partitioned into 28 equal-sized groups of 676 channels each.
Structures are scored and ranked in a unified list across both component types, enabling cross-component pruning decisions: a total of 395{,}232 parameters are removed under both Standard and \rasp{} configurations, ensuring matched parameter budgets.

\paragraph{Importance scoring.}
For each structure $h$, we compute Taylor-expansion importance $I(h)$ (Eq.~\ref{eq:importance}) using activations and gradients collected on 128 randomly sampled chain-of-thought examples from the GSM8K training set.
For GQA groups, we aggregate importance across the constituent query heads and KV head (Section~\ref{sec:stage1}).

\paragraph{CKA-based functional overlap.}
We compute intra-layer linear CKA (Eq.~\ref{eq:cka}) between all pairs of GQA groups and between all pairs of MLP channel groups within each layer, using the same 128 calibration examples.
The functional overlap $F(h)$ for each structure is the mean pairwise CKA with its within-layer peers (Eq.~\ref{eq:overlap}).

\paragraph{Recovery-aware scoring.}
The final pruning score $S(h) = I(h) \times (1 - F(h))^\alpha$ (Eq.~\ref{eq:score}) combines importance and functional overlap with $\alpha{=}1.0$.
Structures with the lowest $S(h)$ are pruned first, subject to the per-layer capacity floor.

\paragraph{Per-layer capacity floor.}
To prevent representational collapse in individual layers, we enforce that each layer retains at least 40\% of its MLP channel groups (${\ge}12$ of 28 groups per layer).
During structure selection, candidates whose removal would violate this constraint are skipped in favor of the next lowest-scoring structure.
This constraint is necessary because unconstrained global ranking at 20\% sparsity concentrates pruning in layers with uniformly high functional overlap, collapsing those layers' capacity.

\paragraph{Baseline: Standard importance-only pruning.}
As the primary baseline, we apply the same Taylor-expansion importance scoring and global ranking but without the CKA overlap penalty ($\alpha{=}0$ in Eq.~\ref{eq:score}), following LLM-Pruner~\citep{llmpruner2023}.
The same per-layer capacity floor and sparsity target are applied, ensuring that the only difference between Standard and \rasp{} is the head selection criterion.

\subsection{Recovery Pipeline}
\label{sec:recovery_pipeline}

After pruning, we apply a two-stage recovery pipeline (Section~\ref{sec:stage3}) consisting of SFT warmup followed by GRPO-based RLVR.

\paragraph{Stage 3a: SFT warmup.}
Structured pruning destroys the model's ability to generate EOS tokens, causing unterminated repetitive outputs that prevent GRPO from computing rewards.
We restore generation competence with a brief SFT phase on 75 chain-of-thought examples (approximately 1\% of the GSM8K training set), trained for 2 epochs (1{,}870 gradient steps) with teacher forcing.
This stage is intentionally minimal: its sole purpose is to restore EOS generation and basic format compliance, not to recover reasoning ability.

\paragraph{Stage 3b: GRPO recovery.}
With generation competence restored, we apply GRPO~\citep{shao2024grpo} on GSM8K training problems to recover reasoning ability.
The reward is binary: $r{=}1$ if the extracted final answer matches the ground truth, $r{=}0$ otherwise.
We use the following configuration: learning rate $1{\times}10^{-6}$, maximum completion length 1{,}024 tokens, 1{,}000 training samples, 1 epoch, KL penalty coefficient $\beta{=}0$, sampling temperature 0.7, and DAPO-style clipped loss.
Setting $\beta{=}0$ removes the KL divergence penalty against the reference policy, allowing the pruned model to freely explore new reasoning paths unconstrained by its pre-pruning behavior---a deliberate choice for post-pruning recovery where the reference policy's reasoning trajectories may no longer be achievable through the reduced architecture.

\subsection{Results}
\label{sec:results}

% [TBD: Results after GRPO completion]
%
% Planned content:
% - Table 1 (\ref{tab:postprune}): Post-prune accuracy comparison
%   Known: RASP 29% vs Standard 1% (N=100, p<0.0001)
% - Table 2 (\ref{tab:recovery}): Full recovery pipeline comparison
%   Known: Post-SFT RASP 57% vs Standard 64% (N=100)
%   Pending: Post-GRPO accuracy for both configurations
% - Analysis of recovery Delta and whether RASP's post-prune advantage
%   translates to superior post-GRPO performance
% - Statistical significance testing (Fisher's exact test, N=1319)

\subsection{Ablation Studies}
\label{sec:ablation}

% [TBD: Ablation results after GRPO completion]
%
% Planned ablations:
% - Table 3 (\ref{tab:alpha}): Effect of CKA balancing coefficient alpha
%   at 20% sparsity (alpha in {0.0, 0.5, 1.0, 2.0, 4.0})
%   alpha=0.0 reduces to Standard importance-only pruning
% - Per-layer capacity floor sensitivity
% - CKA calibration set size (64, 128, 256, 512 samples)
